\section{Total-variation and separation cutoff}\label{sec:count-cutoff}

For a finite chain with semigroup $P_t$ and stationary law $\pi$, write
\[
 d_x(t)=\|P_t(x,\cdot)-\pi\|_{\mathrm{TV}},\qquad
 d(t)=\max_x d_x(t),\qquad
 t_{\mathrm{mix}}(\varepsilon)=\inf\{t:d(t)\le\varepsilon\}.
\]
The superscripts ``count'' and ``polymer'' specify the chain. For a
specified initial state the maximum is omitted. All logarithms are
natural. The time-window statements below imply cutoff because the
window is smaller than the stated mixing time by a factor tending to
zero.

\subsection{The two polymer endpoints}
Let $L=L_M$ be the polymer generator, with stationary law $\pi=\pi_M$,
and write $K(x)$ for the number of polymers. The count strongly lumps
$L$ to the birth--death generator $Q$ on $\{1,\ldots,n\}$. Let $\bar\pi$ denote the pushforward of $\pi$ by $K$.
For every function $g$ on $\{1,\ldots,n\}$,
\[
 L(g\circ K)=(Qg)\circ K,
 \qquad P_t(g\circ K)=(\bar P_tg)\circ K.
\]
Fix $i$ such that the fiber $K^{-1}(i)$ consists of one state $x_i$.
The initial density with respect to stationarity is then
\[
 h_0(x)=\frac{\mathbf1_{\{x_i\}}(x)}{\pi(x_i)}
       =\frac{\mathbf1_{\{i\}}(K(x))}{\bar\pi(i)}.
\]
By reversibility, the density at time $t$ is $P_th_0$; invariance of the
count-function subspace gives
\[
 \frac{P_t(x_i,x)}{\pi(x)}
 =\frac{\bar P_t(i,K(x))}{\bar\pi(K(x))}.
\]
Consequently, at every $t\geq0$,
\[
 \boxed{\ \|P_t(x_i,\cdot)-\pi\|_{\rm TV}
       =\|\bar P_t(i,\cdot)-\bar\pi\|_{\rm TV}.\ }
\]
Indeed, sum the absolute density difference over each count fiber and use
$\sum_{K(x)=k}\pi(x)=\bar\pi(k)$. Both $i=1$ (one polymer of size $n$)
and $i=n$ (all monomers) have singleton fibers, so the identity holds at
both endpoints of the polymer chain. It transfers endpoint cutoff, and
provides a lower bound on worst-case polymer mixing. It does not provide
an upper bound for arbitrary initial polymer states.


\subsection{Count cutoff for fixed scaled fragmentation rate}
Throughout this section, $b>0$ is fixed and constants may depend on $b$.
Let $Q_n$ be the continuous-time birth--death generator on $\{1,\ldots,n\}$
with rates
\[
 \lambda_k=\frac{b(n-k)}n,\qquad \mu_k=\frac{k(k-1)}n.
\]
Its stationary law $p=p_n$ satisfies
\[
 p(k)\propto \binom nk\frac{b^k}{(k-1)!},\qquad
 \frac{p(k+1)}{p(k)}=\frac{b(n-k)}{k(k+1)}.
\]
Define the positive drift zero
\[
 a=a_n=\frac{-(b-1)+\sqrt{(b-1)^2+4bn}}2
       =\sqrt{bn}+O_b(1),
\]
so that
\[
 D_k:=\mu_k-\lambda_k
     =\frac{(k-a)(k+a+b-1)}n.
\]

\begin{theorem}[Count cutoff and polymer endpoint cutoff]\label{thm:count-cutoff}
For every fixed $b>0$, the count chains have a worst-case total-variation
cutoff at
\[
 t_n=\frac{\sqrt n\log n}{8\sqrt b}
\]
with an $O_b(\sqrt n)$ window. Precisely, for each $\varepsilon\in(0,1)$,
\[
 t^{\mathrm{count}}_{\mathrm{mix},n}(\varepsilon)
 =t_n+O_{b,\varepsilon}(\sqrt n).
\]
The same statement holds separately for the count chain started from $1$
and from $n$. The polymer chain at $\beta=b/n$ strongly lumps to $Q_n$. Its
total-variation distance from the single $n$-polymer and all-monomer states
equals the corresponding count distance at every time. These two polymer
initial states therefore have the same cutoff and window. In particular,
the worst-case polymer mixing time is at least
$t_n-O_{b,\varepsilon}(\sqrt n)$. An arbitrary-start upper bound is proved in
Theorem~\ref{thm:full-polymer-tv-upper-new}.
\end{theorem}

\subsection{Elementary stationary estimates}
Write $F_k=p([1,k])$ and $G_k=p([k,n])$. We shall use the following bounds:
\begin{align}
 \mathbb E_p|K-a|&\leq C_b\sqrt a,\label{eq:stationary-first-moment}\\
 c_{b,L}a^{-1/2}\leq p(k)&\leq C_{b,L}a^{-1/2}
       &&(|k-a|\leq L\sqrt a),\label{eq:stationary-local}\\
 F_k\wedge G_k&\geq c_{b,L}
       &&(|k-a|\leq L\sqrt a).\label{eq:stationary-central-masses}
\end{align}
Here $L$ is any fixed finite constant and $n$ is sufficiently large.
In particular every stationary median $m$ satisfies
$|m-a|\leq C_b\sqrt a$.

For completeness, these estimates follow directly from the consecutive
weight ratio. For unnormalized stationary weights $w_k$, set
\[
 R(x)=\frac{b(n-x)}{x(x+1)}.
\]
It is decreasing, and its unique positive solution to $R(x)=1$ is
$a+O_b(1)$. A maximizing index $r$ for $w_k$ therefore satisfies
$r=a+O_b(1)$. On $[a/2,2a]$, for all sufficiently large $n$,
\[
 -\frac{C_b}{a}\leq (\log R)'(x)
 =-\frac1{n-x}-\frac1x-\frac1{x+1}
 \leq-\frac{c_b}{a}.
\]
Since $R(r-1)\geq1\geq R(r)$, summing these derivative bounds over
successive ratios gives
\[
 c_b\exp\!\left(-C_b\frac{(k-a)^2}{a}\right)
 \leq\frac{w_k}{w_r}
 \leq C_b\exp\!\left(-c_b\frac{(k-a)^2}{a}\right),
 \qquad a/2\leq k\leq2a.
\]
Rounding endpoints changes only the constants. Below $a/2$, the successive
ratios in the direction away from the mode are bounded by a fixed number
$q_b<1$; the same holds above $2a$. These two tails are thus geometric,
starting at weights at most $C_be^{-c_ba}w_r$.
Consequently
\[
 \sum_k w_k\asymp_b\sqrt a\,w_r,\qquad
 \sum_k|k-a|w_k\leq C_b a\,w_r.
\]
This proves \eqref{eq:stationary-first-moment} and
\eqref{eq:stationary-local}. To obtain
\eqref{eq:stationary-central-masses}, sum the local lower bound over an
interval of length $\lfloor\sqrt a\rfloor$ immediately to the left, or
right, of $k$, using \eqref{eq:stationary-local} with $L+2$.
Finally, Markov's inequality applied to
\eqref{eq:stationary-first-moment} puts every median within
$C_b\sqrt a$ of $a$ (enlarging $C_b$ as needed).


\subsection{Endpoint hitting-time means}
Let $T_j$ be the first hitting time of $j$ and define the one-edge passage
means
\[
 e_k=\mathbb E_kT_{k+1}=\frac{F_k}{p(k)\lambda_k},\qquad
 f_k=\mathbb E_kT_{k-1}=\frac{G_k}{p(k)\mu_k}.
\]
These standard identities also follow by summing the Poisson equation
$Qh=-1$ against $p$ up to the absorbing edge. The recurrences are
\[
 \lambda_ke_k-\mu_ke_{k-1}=1,\qquad
 \mu_kf_k-\lambda_kf_{k+1}=1,
\]
with $e_1=1/\lambda_1$ and $f_n=1/\mu_n$.
Since $\lambda_k$ decreases and $\mu_k$ increases, induction in the first
recurrence shows that $e_k$ increases with $k$; reverse induction shows
that $f_k$ decreases with $k$.

Fix a stationary median $m$, and choose a fixed $L$ so large that
$|m-a|\leq(L-2)\sqrt a$. Put
$u=\lfloor a-L\sqrt a\rfloor$ and
$v=\lceil a+L\sqrt a\rceil$.
The stationary estimates imply, uniformly on the central interval
$[u,v]$,
\[
 e_k+f_k\leq C_{b,L}\sqrt a
\]
whenever the respective passage is defined: indeed
$p(k)\geq c_{b,L}/\sqrt a$ and
$\lambda_k\wedge\mu_k\geq c_b>0$ there. Hence all passages between
$u$, $m$, and $v$ together contribute $O_{b,L}(a)$ to their mean times.

For $k\leq u$ write $\delta_k=\lambda_k-\mu_k=-D_k>0$ and
$h_k=1/\delta_k$. The recurrence gives $e_k\leq h_k$ by induction, since
$h_k$ increases. If $\Delta_k=h_k-e_k$ and
$q_k=\mu_k/\lambda_k$, then
\[
 \Delta_1=0,\qquad
 \Delta_k=q_k\Delta_{k-1}+q_k(h_k-h_{k-1}).
\]
The sequence $q_k$ increases, and
\[
 h_k-h_{k-1}
 =\frac{2k+b-2}{n\delta_k\delta_{k-1}}
\]
also increases. Iterating and bounding the resulting geometric series
therefore gives
\[
 0\leq h_k-e_k
 \leq\frac{q_k}{1-q_k}(h_k-h_{k-1})
 \leq\frac{\mu_k(\delta_{k-1}-\delta_k)}{\delta_k^3}
 \leq\frac{C_ba^2}{(a-k)^3}.
\]
The last inequality uses $n\asymp_ba^2$, $\mu_k\leq C_b$,
$\delta_{k-1}-\delta_k\leq C_b/a$, and
$\delta_k\geq c_b(a-k)/a$. Summing for $k\leq u$ yields
\[
 \sum_{k=1}^{u}e_k
 =\sum_{k=1}^{u}\frac1{\delta_k}+O_{b,L}(a),
\]
because $\sum_{d\geq L\sqrt a-1}d^{-3}=O_L(a^{-1})$.

For $k\geq v$, put $h_k=1/D_k$. This time $h_k$ decreases, and
$h_k-h_{k+1}$ decreases as well: the function
$x\mapsto n/[x^2+(b-1)x-bn]$ is decreasing and convex for $x>a$
(and $n$ large). The ratio $p_k^*=\lambda_k/\mu_k$ decreases. Reverse
iteration of the recurrence, starting at $f_n=h_n$, gives
\[
 0\leq h_k-f_k
 \leq\frac{p_k^*}{1-p_k^*}(h_k-h_{k+1})
 \leq\frac{\lambda_k(D_{k+1}-D_k)}{D_k^3}
 \leq\frac{C_ba^2}{(k-a)^3}.
\]
At $k=n$ the difference is zero. For the final estimate, use
$\lambda_k\leq b$, $D_{k+1}-D_k=(2k+b)/n$, and
$D_k\geq c_b(k-a)(k+a)/n$; hence the displayed ratio is at most
$C_bn^2/[(k-a)^3(k+a)^2]\leq C_ba^2/(k-a)^3$.
It follows that
\[
 \sum_{k=v}^{n}f_k
 =\sum_{k=v}^{n}\frac1{D_k}+O_{b,L}(a).
\]

The reciprocal drift sums have elementary partial fractions:
\begin{align*}
 \sum_{k=1}^{u}\frac1{\delta_k}
 &=\frac n{2a+b-1}
   \sum_{k=1}^{u}\left(\frac1{a-k}+\frac1{a+k+b-1}\right),\\
 \sum_{k=v}^{n}\frac1{D_k}
 &=\frac n{2a+b-1}
   \sum_{k=v}^{n}\left(\frac1{k-a}-\frac1{k+a+b-1}\right).
\end{align*}
Integral comparison for the decreasing functions $1/x$ shows that each
bracketed sum is
$\tfrac12\log a+O_{b,L}(1)$: the first singular sum runs from a distance
of order $a$ down to $L\sqrt a$, and in the second sum the upper endpoint
contributes only
$\log[(n-a)/(n+a+b-1)]=O_b(a/n)$ after cancellation.
Thus, by the strong Markov property and addition of passage means,
\[
 \boxed{\begin{aligned}
 \mathbb E_iT_m
 &=\frac n{4a}\log a+O_b(a)\\
 &=\frac{\sqrt n\log n}{8\sqrt b}+O_b(\sqrt n),
 \qquad i\in\{1,n\}.
 \end{aligned}}
\]
More explicitly, both endpoint means separately equal the last expression.
The same proof works for any target $j$ with $|j-a|\leq L_0\sqrt a$,
with constants allowed to depend on the fixed $L_0$.


\subsection{Endpoint hitting-time variances}
Successive one-edge passage times are independent by the strong Markov
property. We prove
\[
 \operatorname{Var}_1T_m+\operatorname{Var}_nT_m=O_b(a^2)=O_b(n).
\]
For an upward passage $T=T_{k+1}$ started at $k$, let
$H(j)=\mathbb E_jT_{k+1}=\sum_{\ell=j}^{k}e_\ell$ for $j\leq k$.
The equations $QH=-1$ and
$Q\mathbb E_{\cdot}T_{k+1}^2=-2H$ on $\{1,\ldots,k\}$,
with zero boundary value at $k+1$, give upon summing against $p$
\begin{align*}
 \mathbb E_kT_{k+1}^2
 &=\frac2{p(k)\lambda_k}\sum_{j=1}^{k}p(j)H(j)\\
 &=\frac2{p(k)\lambda_k}\sum_{\ell=1}^{k}e_\ell F_\ell
 \leq2e_k^2\sum_{\ell=1}^{k}\frac{F_\ell}{F_k}.
\end{align*}
The inequality uses monotonicity of $e_\ell$ and
$e_k=F_k/[p(k)\lambda_k]$. In particular,
\begin{equation}\label{eq:upward-variance-bound}
 \operatorname{Var}_kT_{k+1}
 \leq2e_k^2\sum_{\ell=1}^{k}\frac{F_\ell}{F_k}.
\end{equation}
The reflected argument, using monotonicity of $f_k$, gives
\begin{equation}\label{eq:downward-variance-bound}
 \operatorname{Var}_kT_{k-1}
 \leq2f_k^2\sum_{\ell=k}^{n}\frac{G_\ell}{G_k}.
\end{equation}

For $k\leq u=\lfloor a-L\sqrt a\rfloor$, set
\[
 r_k=\frac{p(k-1)}{p(k)}
     =\frac{\mu_k}{\lambda_{k-1}}
     \leq\frac{\mu_k}{\lambda_k}<1,
\]
with $r_1=0$. These ratios increase with $k$. Shifting every index in
the sum defining $F_\ell$ by $k-\ell$ gives
$F_\ell\leq r_k^{k-\ell}F_k$. Consequently,
\[
 \sum_{\ell=1}^{k}\frac{F_\ell}{F_k}
 \leq\frac1{1-r_k}\leq\frac{\lambda_k}{\delta_k}.
\]
Together with $e_k\leq1/\delta_k$, this implies
\[
 \operatorname{Var}_kT_{k+1}
 \leq\frac{2\lambda_k}{\delta_k^3}
 \leq\frac{C_ba^3}{(a-k)^3}.
\]
Summing over $k\leq u$ gives $O_b(a^2)$.

For $k\geq v=\lceil a+L\sqrt a\rceil$, use instead the decreasing
ratios
\[
 s_k=\frac{p(k+1)}{p(k)}
     =\frac{\lambda_k}{\mu_{k+1}}
     \leq\frac{\lambda_k}{\mu_k}<1,
\]
with $s_n=0$. The same index-shifting argument gives
\[
 \sum_{\ell=k}^{n}\frac{G_\ell}{G_k}
 \leq\frac1{1-s_k}\leq\frac{\mu_k}{D_k}.
\]
Since $f_k\leq1/D_k$,
\[
 \operatorname{Var}_kT_{k-1}
 \leq\frac{2\mu_k}{D_k^3}.
\]
For $v\leq k\leq2a$, this is at most
$C_ba^3/(k-a)^3$, whose sum is $O_b(a^2)$.
For $k>2a$, use $D_k\geq c_bk^2/n$ and $\mu_k\leq k^2/n$
to obtain the bound $C_bn^2/k^4$; its sum is
$O_b(n^2/a^3)=O_b(a)$.

It remains to treat the $O_b(\sqrt a)$ central passages. Uniformly for
$k\in[u,v]$, the stationary estimates and the mean formulas give
$e_k+f_k\leq C_b\sqrt a$ and $F_k\wedge G_k\geq c_b>0$.
Furthermore,
\begin{align*}
 \sum_{\ell=1}^{k}F_\ell
 &=\mathbb E_p[(k-K+1)_+]\leq|k-a|+1+\mathbb E_p|K-a|
   \leq C_b\sqrt a,\\
 \sum_{\ell=k}^{n}G_\ell
 &=\mathbb E_p[(K-k+1)_+]\leq C_b\sqrt a.
\end{align*}
Equations \eqref{eq:upward-variance-bound} and
\eqref{eq:downward-variance-bound} therefore bound each central passage
variance by $C_ba^{3/2}$. Their total is $O_b(a^2)$.
Adding independent passage variances proves
\[
 \boxed{\quad
 \operatorname{Var}_1T_m\vee\operatorname{Var}_nT_m
 \leq C_ba^2\leq C_bn.
 \quad}
\]
As with the means, this bound remains valid for every target within a
fixed multiple of $\sqrt a$ of $a$, with the constant depending on that
multiple.


\subsection{Application of the birth--death cutoff theorem}
We use Chen--Saloff-Coste \cite[Theorem~1.4(1)]{CSC2015}.
In the continuous-time setting, if $m_n$ has
stationary mass bounded below on both sides, put
\[
 \theta_n=\max\{\mathbb E_1T_{m_n},\mathbb E_nT_{m_n}\},\qquad
 \alpha_n^2=\max\{\operatorname{Var}_1T_{m_n},
                         \operatorname{Var}_nT_{m_n}\}.
\]
Their theorem states that $\theta_n/\alpha_n\to\infty$ implies a
worst-case total-variation $(\theta_n,\alpha_n)$ cutoff. Its published
normalization starts from a stochastic matrix $B_n$ and uses the generator
$B_n-I$. To apply it here, take
\[
 c_n\geq\max_k(\lambda_k+\mu_k),\qquad B_n=I+Q_n/c_n.
\]
The continuous-time chain with generator $B_n-I$ is our chain slowed by
$c_n$. All passage means and standard deviations are multiplied by $c_n$;
both the criterion and its time-window conclusion therefore transfer
exactly to $Q_n$.

Take $m_n$ to be any stationary median. The previous two subsections give
\[
 \theta_n=t_n+O_b(a),\qquad \alpha_n=O_b(a),\qquad
 \frac{\theta_n}{\alpha_n}\longrightarrow\infty.
\]
The theorem yields
\[
 t^{\mathrm{count}}_{\mathrm{mix},n}(\varepsilon)
 =t_n+O_{b,\varepsilon}(a)
 =t_n+O_{b,\varepsilon}(\sqrt n).
\]
This proves the worst-case count assertion of
Theorem~\ref{thm:count-cutoff}.

To prove the assertion for each endpoint separately, it remains only to
give endpoint lower bounds; their upper bounds follow from the worst-case
bound. For a fixed $L>0$, let
$j_-=\lfloor a-L\sqrt a\rfloor$ and
$j_+=\lceil a+L\sqrt a\rceil$.
The stationary moment estimate gives
\[
 p([1,j_--1])\vee p([j_++1,n])\leq C_b/L.
\]
The preceding passage estimates, now with these central targets, give
\[
 \left|\mathbb E_1T_{j_-}-t_n\right|
 +\left|\mathbb E_nT_{j_+}-t_n\right|\leq C_{b,L}a,
 \qquad
 \operatorname{Var}_1T_{j_-}\vee\operatorname{Var}_nT_{j_+}
 \leq C_{b,L}a^2.
\]
By nearest-neighbor paths,
$\{T_{j_-}>t\}\subseteq\{K_t<j_-\}$ when $K_0=1$, and
$\{T_{j_+}>t\}\subseteq\{K_t>j_+\}$ when $K_0=n$.
Chebyshev's inequality therefore gives, for $A>C_{b,L}$ and $n$ large,
\[
 d_{n,1}(t_n-Aa)\wedge d_{n,n}(t_n-Aa)
 \geq1-\frac{C_b}{L}-\frac{C_{b,L}}{(A-C_{b,L})^2}.
\]
First choose $L$ and then $A$ large. Together with the count upper bound,
this proves both endpoint cutoffs with an $O_b(\sqrt n)$ window. The exact
endpoint identity in the first subsection completes the polymer endpoint
claims.

\subsection{A further consequence: separation cutoff}
Define separation from an initial state $x$ by
\[
 \operatorname{sep}_x(t)=1-\min_y\frac{P_t(x,y)}{\pi(y)}.
\]
The same established passage estimates give
\[
 s_n:=\mathbb E_1T_m+\mathbb E_nT_m
      =\frac{\sqrt n\log n}{4\sqrt b}+O_b(\sqrt n),\qquad
 v_n:=\sqrt{\operatorname{Var}_1T_m+\operatorname{Var}_nT_m}
      =O_b(\sqrt n).
\]
Chen--Saloff-Coste \cite[Theorem~1.3(1)]{CSC2015} implies endpoint separation cutoff at
$s_n$ with window $v_n$. In a continuous-time birth--death chain maximum
separation is attained from either boundary state; see their
Remark~1.2. Thus
\[
 t^{\mathrm{count}}_{\mathrm{sep},n}(\varepsilon)
 =\frac{\sqrt n\log n}{4\sqrt b}
       +O_{b,\varepsilon}(\sqrt n).
\]
The singleton-fiber identity gives a density that is constant on every
count fiber. Taking the minimum of that density, rather than summing its
absolute difference from one, proves the exact endpoint separation
identity for the polymer chain as well. Both polymer endpoint separation
mixing times therefore have this same asymptotic formula. This again does
not yield an arbitrary-start polymer upper bound.
